\cotefile{157-1}
\scanpage{50}
\mspage{56}

Cela provient du fait que pour un tel foncteur \unsure{propr.|propre}, \deleted{pour tout faisceau $F$ \unsure{abélien|d'ens}} si $F$ est un préfaisceau \unsure{ou|en} complexe de préfaisceaux abéliens \inserted{sur $X$}, la formation de $Rf_* (F)$ commute au passage aux fibres, et de $=$ pour $R^1f_*(F)$ si $F$ est un \deleted{faisceau}préfaisceau en groupes (pas néc abéliens).

\uline{Cor.} 1 Si $f:X\longrightarrow Y$ est coh. propre à fibres contractiles, alors $Y$ contractile \deleted{\unsure{alors|donc}} $X$ l'est.

Grâce à ceci, on trouve que le composé de deux foncteurs propres (lisses) est propre (lisse).

\uline{NB} Contrairement à ce qui a lieu en topologie pour la notion de propreté, si $gf$ est propre et $g$ propre, $f$ ne l'est pas néc.$^{\mathrm{t}}$, \deleted{\unsure{p. ex. les}} \inserted{i.e un} $S$-morphisme de $S$-cat $X,Y$ propres/$S$ n'est pas néc. propre. Ainsi une section de $X/S$, pour $X$ propre sur $S$, n'est pas néc. propre i.e. n'est pas néc. une immersion fermée. C'est déjà faux si $S=e$ (auquel cas $f:X\longrightarrow e$ est tjrs propre).

\uline{Cor 2} \deleted{\unsure{Un}} morphisme propre ou lisse, à fibres contractiles, est à la fois une hot équivalence stricte, et costricte. \unsure{En effet}, \unsure{si $X_y$|$X_y$} \unsure{$Y_{/y}$|$Y_y$} \unsure{car|ce} \unsure{(fibr.)|Cofibr.} \unsure{et} $X_{/y}\longrightarrow Y_{/y}$

\uline{Théorème} (g) Soit une \unsure{situation} de \deleted{\unsure{de}} chang$^{\mathrm{mt}}$ de base \unsure{dans|du} \unsure{Diag|Dér.|Dév.}

\diagram{\begin{tikzcd}[ampersand replacement=\&]
X \arrow[d,"f"'] \& X' \arrow[l,leftrightarrow,"q"'] \arrow[d,"f'"] \\
Y \& Y' \arrow[l,"p"] \arrow[ul,phantom,"\mathrm{cart}" description]
\end{tikzcd}}

\[
,\qquad p^*f_*(F)\xrightarrow{c}f'_*q^*(F).
\]

($F\in\unsure{\mathbb D|\underline D}(X)$, $\underline D$ dérivateur \unsure{assoc.|satisf.} à \unsure{gr.|pr.} Si $f$ est (homot.) propre, \uline{ou} $p$ hom$^{\mathrm{t}}$ lisse, \deleted{\unsure{alors}} $c$ est iso.
% archive sheet number: 49
\scanpage{51}
\mspage{57}

\unsure{La} question, c'est si pour tout $y'\in Y'$, posant $y=p(y')$, \deleted{\unsure{et}} \deleted{\unsure{cons}}considérant le foncteur induit par $q$

\diagram{\begin{tikzcd}[ampersand replacement=\&]
X_{/y} \& X'_{/y'} \arrow[l]
\end{tikzcd}}

le morphisme can.

\[
\Gamma_{\unsure{\mathbb D|\underline D}}(X_{/y},F)\unsure{\relbar\joinrel\relbar|\longrightarrow}\Gamma_{\mathbb D}(X'_{/y'},q^*F)
\]

est iso. Dans le cas $f$ propre, on considère

\diagram{\begin{tikzcd}[ampersand replacement=\&]
X_y \arrow[d,"{\sim\ \text{costricte}}"'] \& X_{y'} \arrow[l,"\sim"'] \arrow[d,"{\sim\ \text{costricte}}"] \\
X_{/y} \& X'_{/y'} \arrow[l]
\end{tikzcd}}

d'où la conclusion puisque par \unsure{Dèr|Dér} 5 \unsure{g|G}, dans

\diagram{\begin{tikzcd}[ampersand replacement=\&]
\Gamma(X_y) \arrow[r] \& \Gamma(X_{y'}) \\
\Gamma(X_{/y}) \arrow[u] \arrow[r] \& \Gamma(X'_{/y'}) \arrow[u]
\end{tikzcd}}

les flèches verticales et l'hor. sup. sont iso -- donc aussi l'hor. inf. \deleted{\unsure{Dans}}

Dans le cas $p$ lisse, considérons du côté $Y/Y'$ le carré de chang$^{\mathrm{ts}}$ de base

\diagram{\begin{tikzcd}[ampersand replacement=\&]
Y \& Y' \arrow[l] \\
Y_{/y} \arrow[u] \& Y'_{/y'} \arrow[u] \arrow[l]
\end{tikzcd}}
% archive sheet number: 50
\scanpage{52}
\mspage{58}

d'où le diagramme cubique

\diagram{\begin{tikzcd}[ampersand replacement=\&,column sep=small,row sep=small]
X \arrow[dd] \& {} \& X' \arrow[ll] \arrow[dd] \& {} \\
{} \& X_{/y} \arrow[ul] \& {} \& X'_{/y'} \arrow[ul] \arrow[ll,"\unsure{k|\kappa}"'] \arrow[dd] \\
Y \& {} \& Y' \arrow[ll] \& {} \\
{} \& Y_{/y} \arrow[ul] \& {} \& Y'_{/y'} \arrow[ul] \arrow[ll,"\pi"] \arrow[from=2-2,to=4-2,crossing over]
\end{tikzcd}}

, où les \deleted{\unsure{côtés|cotés}} carrés latéraux (verticaux) sont cartésiens. On note que $\pi$ est lisse et à fibres \deleted{\unsure{car}}contractiles (\unsure{car|si} $p$ est lisse) donc $k$ aussi. \deleted{Notons \unsure{donc|alors}} Donc $k$ est une hot-équivalence \deleted{\unsure{stricte}} \deleted{\unsure{Cor.}} \deleted{\unsure{les}} \unsure{$=$} costricte \unsure{$g$|\}}, donc $\Gamma(X_{/y})\mathrel{\overset{\sim}{\unsure{\relbar\joinrel\relbar|\longrightarrow}}}\Gamma(X'_{/y'})$, on gagne encore.

\uline{NB} Sauf erreur, on a une réciproque du th., savoir : Si $f:X\longrightarrow Y$ est tel que le th. de chgt. de base pour $f_*(F)$ est vrai pour \uline{tout} dérivateur (il suffit sans doute de $\mathrm{HOTAB}$ et $\mathrm{HOT}$) et pour tout $Y'\unsure{\relbar\joinrel\relbar|\longrightarrow}Y$, alors $f$ est hom$^{\unsure{\mathrm t}}$ propre.

D'autre part, si $g:Y'\longrightarrow Y$ est tel que le th. de chang$^{\mathrm t}$ de base par $g$ est vrai pour tout $X\xrightarrow{f}Y$ et tout $F\in\unsure{\mathbb D|\underline D}(X)$ ($\underline D$ dérivateur \inserted{\unsure{pr. à dr.|gr. à dr.}} quelc.), alors $g$ est hom$^{\mathrm t}$ lisse.

Du côté homologique : les foncteurs \uline{lisses} $f$ sont caractérisés comme ceux pour lesquels \deleted{\unsure{R}} $f_!$ commute à tout changement de base. Les foncteurs propres $g$ sont exactement ceux pour lesquels le chang$^{\mathrm t}$ de base par $g$ commute à $f_!$, pour \deleted{\unsure{ex. :}} \inserted{tout} $f:X\longrightarrow Y$ et tout $\underline D$ \unsure{dér}
% archive sheet number: 51
\scanpage{53}
\mspage{59}

Ainsi, les foncteurs propres $f:X\longrightarrow Y$ se caractérisent de \uline{quatre} façons remarquables

a) $f_*$ commute à tout changement de base ($\forall$ dérivateur \inserted{$\unsure{\mathbb D|\underline D}$} exact à g.)

b) le chang$^{\mathrm{mt}}$ de base $f^*$ commute à tout $g_!$ ($g:\deleted{\unsure{X}}Y'\unsure{\relbar\joinrel\relbar|\longrightarrow}Y$, $\mathbb D$ dér. exact à droite)

c) $\forall\deleted{\unsure{x}}y\in\mathrm{Ob}\,\deleted{\unsure{X}}Y$, $\underset{\text{fibre de }X\text{ en }y}{X_{\deleted{\unsure{x}}y}}\longrightarrow X_{/y\deleted{\unsure{f(x)}}}$ est une hot. équiv. costricte \inserted{i.e. est un foncteur homotop$^{\mathrm t}$ cofinal}

d) $\forall x\in\mathrm{Ob}\,X$, \deleted{\unsure{$x\mapsto$}} le foncteur induit sur les colocalisés

\[
{}_{x\backslash}X_{\deleted{\unsure{/x}}}\longrightarrow{}_{f(x)\backslash}Y_{\deleted{\unsure{/f(x)}}},
\]

est à fibres contractiles.

\[ \rule{1.3cm}{0.4pt} \]

Autre situation \unsure{standart|standard} où le chang$^{\mathrm t}$ de base marche (mais où les foncteurs qui interviennent n'ont aucune raison d'être propres ou \unsure{lisses|lisse}).

Soient

\[ f:I\longrightarrow I',\qquad g:J\longrightarrow J' \]

de $\mathrm{Diag}$, d'où le diagramme (comparer page 33)

\diagram{\begin{tikzcd}[ampersand replacement=\&]
X'=I\times J \arrow[r,"G'"] \arrow[d,"{\deleted{\unsure{G}}F'}"'] \& I\times J'\qquad\unsure{X|X'} \arrow[d,"F"] \\
Y'=I'\times J \arrow[r,"G"] \& I'\times J'=Y
\end{tikzcd}}

Ce diagramme \deleted{\unsure{satis}} est tel que pour tt dér. à g. $\mathbb D$,
% archive sheet number: 52
\scanpage{54}
\mspage{60}

on \unsure{a} le changement de base

\[ G^*F_*\simeq F'_*G'^{\unsure{*|\wedge}}, \]

la raison ici étant que les foncteurs induits

\diagram{\begin{tikzcd}[ampersand replacement=\&]
X'_{/y'} \arrow[r,phantom,"{\unsure{\relbar\joinrel\relbar|\longrightarrow}}" description] \& X_{/y}
\end{tikzcd}}

sont des hot-équiv. \deleted{\unsure{non}} \inserted{co}strictes, i.e. que les catégories $K$ (page \unsure{49|99}) sont contractiles.

En effet, si

\[
\underset{\substack{\Vert\\(i_0,j'_0)}}{x_0}\in\mathrm{Ob}\,X,\quad
\underset{\substack{\Vert\\(i'_0,j_0)}}{y'_0}\in\mathrm{Ob}\,Y',\quad
\underset{\substack{\Vert\\(u'_0,v'_0)}}{u_0}:F(x)\longrightarrow G(y')
\]

\inserted{où $u'_0:f(i_0)\longrightarrow i'_0$ \quad $v'_0:j'_0\longrightarrow g(j_0)$}

sont \unsure{donnés|données}, on vérifie que la catégorie $K$ cherchée est \deleted{\unsure{équivalente}} isom. simplement à

\[
{}_{(\unsure{i_0|i'_0},u'_0)\backslash}\bigl(I_{/i'_0}\bigr)\times{}_{(j_0,v'_0)\backslash}\bigl(J_{/\unsure{j'_0|j_0}}\bigr),
\]

qui est contractile (ayant un objet \deleted{\unsure{final})} initial. \qquad Ouf !

\[ \rule{1.3cm}{0.4pt} \]

Je m'aperçois d'un très gros \unsure{canular}. Je croyais me rappeler que pour $f:X\longrightarrow Y$ qui est une hot-équivalence stricte, on avait pour tout $F\in\unsure{\mathbb D|\underline D}(Y)$ $F\xrightarrow{\sim}f_*f^*F$ (p.ex. pour $\mathbb D=\deleted{\unsure{\mathrm{HOTAB}}}\mathrm{Der}(\mathrm{Ab})$, ou $\mathbb D=\mathrm{HOT}$), i.e. $f^*$ pleinement fidèle, \deleted{\unsure{tout objet} $F$} d'où aussi que $\Gamma_{\mathbb D}(Y,F)\mathrel{\overset{\deleted{\unsure{f^*}}}{\unsure{\relbar\joinrel\relbar|\longrightarrow}}}\Gamma_{\mathbb D}(X,f^*F)$ est iso. Mais \deleted{\unsure{ce}} ce \deleted{\unsure{déjà}} n'est vrai que si $F$ est un coefficient \uline{constant} sur $Y$, i.e. de la forme $q^*(M)$ ($q:Y\longrightarrow e$, $M\in\mathbb D(e)=\underline{A_{\mathbb D}}$).

C'est \deleted{\unsure{faux}} pour $F$ général, déjà dans le cas de l'inclusion $\{\unsure{\varphi|e|y}\}\xrightarrow{f}Y$, où $\varphi$ est un objet initial. \deleted{\unsure{qui est} $f$ \unsure{pourtant}} \inserted{Alors $f$ est bien} une hot-équiv. stricte, mais \deleted{\unsure{$f^*$ n'est pl. fid.}} l'\unsure{hom} $\Gamma(Y,F)\unsure{\relbar\joinrel\relbar|\longrightarrow}F(x)$ n'est \deleted{\unsure{un iso que si}} en général pas un iso.
% archive sheet number: 53
\scanpage{55}
\mspage{61}

\deleted{\unsure{B)}} \uline{Notions de} $\unsure{\mathbb D|\underline D}$-\uline{équivalence} et \uline{de} $\mathbb D$-\uline{contractilité}

On va donc reprendre les questions d'acyclicité et de chgt. de base.

\unsure{Déf.} 1 Soit $\mathbb D$ un \inserted{pré}dérivateur, et soit

\diagram{\begin{tikzcd}[ampersand replacement=\&,row sep=small,column sep=small]
f:X \arrow[rr] \arrow[dr,dotted,"p"'] \& {} \& Y \arrow[dl,dotted,"q"] \\
{} \& e \& {}
\end{tikzcd}}

dans $\mathrm{Diag}$. On dit que $f$ est une \uline{$\mathbb D$-équivalence} si le foncteur

\[ F\longrightarrow f^*F \qquad \mathbb D(Y)\longrightarrow\mathbb D(X) \]

\inserted{induit un foncteur} \deleted{\unsure{est}} pleinement fidèle sur la sous-catégorie de $\mathbb D(Y)$ formée des $\mathbb D$-coefficients \uline{constants}, i.e. si pour tt $M,\inserted{N}\in\mathrm{Ob}(A=\mathbb D(e))$ (``coefficients absolus'') on a

\[ \mathrm{Hom}(q^*M,q^*N)\xrightarrow{\sim}\mathrm{Hom}(p^*M,p^*N) \tag*{(1)} \]

Supposons que $p_*$ et $q_*$ existent. On posera, pour un faisceau $F$ sur $X$

\[ H^{\bullet}_{\mathbb D}(X,F)\overset{\text{\unsure{def|déf}}}{=}p_*(F)\qquad
(\text{et de }{=} \ H_{\bullet}^{\mathbb D}(X,\unsure{P|F})=p_!(F)). \]

\margin{\uline{NB} et si $F=p^*M$, $M\in A$, on note simplement $H^\bullet_{\mathbb D}(X,N)$, $H_\bullet^{\mathbb D}(X,N)$}

L'hom (1) s'interprète \inserted{par adj.} comme

\diagram{\begin{tikzcd}[ampersand replacement=\&,column sep=small]
\mathrm{Hom}(M,q_*q^*N) \arrow[r] \arrow[d,equal,"{\deleted{\unsure{\sim}}}"'] \& \mathrm{Hom}(M,p_*p^*N) \arrow[d,equal,"{\deleted{\unsure{\sim}}}"] \\
\mathrm{Hom}(M,H^\bullet_{\mathbb D}(Y,q^*N)) \& \mathrm{Hom}(N,H^\bullet_{\mathbb D}(X,p^*N)) \arrow[l]
\end{tikzcd}}

déduit de l'hom

\[ q_*q^*N\longrightarrow p_*p^*N\qquad\text{i.e. }H^\bullet_{\mathbb D}(Y,N)\longrightarrow H^\bullet_{\mathbb D}(X,N)\tag*{(2)} \]

Donc dire que (1) est iso pour ts $M,N$ signifie que (2) est un iso pour tout $N$. De $=$, si $p_!,q_!$ existent, cela équivaut à dire que

\[ p_{\unsure{!|*}}p^*N.\longrightarrow q_!q^*N)\quad\text{i.e. }H_\bullet^{\mathbb D}(X,N)\longrightarrow H_\bullet^{\mathbb D}(Y,N)\tag*{($2'$)} \]

est \deleted{\unsure{un}} iso \inserted{dans $A$} pour tout $N\in A$
% archive sheet number: 54
\scanpage{56}
\mspage{62}

En résumé

% margin note written sideways
\margin{\deleted{Si $f_*$ existe, $f$ est une $\unsure{\mathbb D|\underline D}$-équiv. ssi, pour tout faisceau \uuline{constant} $F$ sur $Y$, $F\xrightarrow{\sim} f_*f^*F$ (adj.) est iso. Si $f_!$ existe, $f$ est une $\mathbb D$-équiv. ssi $\forall F$ \uuline{constant} $f_!f^*F\xrightarrow{\sim}F$ (adj) est iso.}}

\uline{Prop. 1} Si \deleted{\unsure{$H^\bullet_{\mathbb D}$}} $p_*$ et $q_*$ existent i.e. si $H^\bullet_{\mathbb D}(X,-)$ et $H^\bullet_{\mathbb D}(Y,\unsure{\cdot|-})$ existent (il suffit $=$ que ça existe pour des coeff. constants) alors $f$ est une $\mathbb D$-équiv. ssi \deleted{\unsure{$\forall N$}} $f$ induit un iso. pour \unsure{la|les} \deleted{\unsure{D}}$\mathbb D$-coh. à coeffts constants de $X$ et $Y$. Si $p_!$ et $q_!$ existent, \inserted{i.e $H_\bullet^{\mathbb D}(X,-)$ et $H_\bullet^{\mathbb D}(Y,-)$ existent} $f$ est une $\mathbb D$-équiv. ssi $f$ induit un iso pour les $\mathbb D$-hom. à coefficients constants de $X$ et $Y$.

\uline{Prop 2} : \deleted{\illegible{1}} les \deleted{\unsure{équivalences}} $\mathbb D$-équiv. Si $X\xrightarrow{f}Y\xrightarrow{g}Z$, si deux parmi $f,g,gf$ sont des $\mathbb D$-équiv, le troisième aussi

\uline{Déf. 2} On dit que $X$ \deleted{\unsure{$\in\mathrm{Ob}$} de $\mathrm{Diag}$} est $\mathbb D$-asphérique (ou $\mathbb D$-contractile, ou $\mathbb D$-triviale) si $X\xrightarrow p e$ est une $\mathbb D$-équivalence, i.e. si $p^*$ est pl. fid.

\deleted{\unsure{Si}} Si $p_*$ existe, cela signifie donc que pour $M\in\mathrm{Ob}\,A$

\[ M\xrightarrow{\overset{\mathrm{adj}}{\sim}}p_*p^*M\qquad\text{i.e. }M\xrightarrow{\sim}H^\bullet_{\mathbb D}(X,M), \]

et si $p_!$ existe, cela signifie que pour tout $M$

\[ p_!p^*M\xrightarrow{\overset{\mathrm{adj}}{\sim}}M\qquad\text{i.e. }H_\bullet^{\mathbb D}(X,M)\xrightarrow{\sim}M. \]

\inserted{Cor.} \deleted{\unsure{NB}} Si $Y$ \deleted{\unsure{$=$}} (resp $X$) est contractile, alors $f:X\longrightarrow Y$ est une $\mathbb D$-équiv. ssi $X$ (resp $Y$) est contractile

\deleted{\uline{\unsure{Déf 3}}}

\uline{Exemple} Prenons $\mathbb D=\mathrm{HOT}$. \inserted{(On peut prendre $\mathrm{Diag}=$ cat des petites catégories)} Alors on a pour $M\in A=\mathrm{Hot}$ et $X$ petite catégorie, \deleted{$p:X\longrightarrow e$}, on a

\[ H^\bullet_{\mathrm{HOT}}(X,M)\simeq\mathrm{Hom}_A(\mathrm{hot}(X),M) \]

où

\[ \mathrm{hot}:\unsure{\{\mathrm{Cat}\}|\mathrm{(Cat)}}\longrightarrow A=\mathrm{Hot} \]

est le foncteur \deleted{qui} can. qui associe à tt $X$ son ``type d'homotopie''. On a

\[ H_\bullet^{\mathrm{HOT}}(X,M)\unsure{=|\simeq}\mathrm{hot}(X)\times M \]

\margin{$\uparrow$ 1} \uuline{Remarques} 1 Pour que $f$ soit $\mathbb D$-équiv, il \uuline{suffit} que pour tt $F$ \uuline{constant} sur $Y$ \inserted{les morphismes d'adjonction} $F\longrightarrow f_*f^*F$ soit iso. \uline{Rem 2} Si $f,f':X\longrightarrow Y$ sont dans une $=$ \unsure{comp.|compos.} connexe de $\uline{\mathrm{Hom}}(X,Y)$, alors $f$ est $\mathbb D$-équiv, $\Leftrightarrow f'$ $\mathbb D$-équiv.
% archive sheet number: 55
\scanpage{57}
\mspage{63}

On en déduit que \deleted{\unsure{la}} \inserted{$f:X\longrightarrow Y$ est une} $\mathrm{HOT}$-équiv. (au sens de la def. 1) ssi c'est une \inserted{hot-}équivalence (ou ``équivalence faible'') i.e. ssi

\[ \mathrm{hot}(f):\mathrm{hot}(X)\xrightarrow{\sim}\mathrm{hot}(Y) \]

est un iso.

\deleted{\unsure{B)}} \uline{Proposition 2} Soit $f:X^{\deleted{\unsure{\prime}}}\xrightarrow{f}\deleted{\unsure{X}}Y$. Conditions équivalentes (si $\unsure{\mathbb D|\underline D}$ dérivateur \uline{à gauche})

a) $f$ est une $\mathbb D$-équiv., et le reste \deleted{pour} \inserted{après} tt chgt. de base $Y'\longrightarrow Y$ qui est un \deleted{\unsure{immersion}} \inserted{isom. local.} \inserted{(p.ex. pour $Y_{/F}\longrightarrow Y$, $F\in\mathrm{Ob}\,Y^{\unsure{\wedge|\prime}}$}

% slanted
\margin{Autre condition équivalente d) Pour tt faisceau \uline{constant} $M_Y$ sur $Y$, on a $M_Y\xrightarrow{\sim}f_*f^*M_Y$ d$'$) Comme d), avec faisceau \uline{loc} constant}

b) \deleted{$f$ \unsure{est une} $\mathbb D$-\unsure{équiv.} $f$} Pour tout $y\in\mathrm{Ob}\,Y$, le foncteur induit $X_{/y}\longrightarrow Y_{/y}$ est une $\mathbb D$-équiv.

c) Pour tout $y\in\mathrm{Ob}\,Y$, $X_{/y}$ est $\mathbb D$-contractile.

\deleted{\unsure{Défin}} N.B $\rightarrow$ On dit que $Z\xrightarrow{g}Y$ est un \uline{iso. local} si $\forall z\in\mathrm{Ob}\,Z$, posant $y=g(z)$, le morphisme induit pour les \deleted{\unsure{cat}} catégories localisées en $z$ et $y$

\[ Z_{/z}\xrightarrow{\sim}Y_{/y} \]

est un iso. Par exemple toute immersion ouverte, tout \deleted{\unsure{morphisme}} \inserted{foncteur} induit $Y_{/y}\longrightarrow Y$ (foncteur de localisation) \deleted{\unsure{sont}} des iso. loc.

Il en résulte a) $\Rightarrow$ b) trivial$^{\mathrm{mt}}$.

b) $\Leftrightarrow$ c). Résulte du cor. de la page précédente, si on sait que \fbox{$Y_{/y}$ est $\mathbb D$-contractile}. Cela \unsure{résultera|résulte} \deleted{de} Der 5 ci-dessous

\margin{$\searrow$ Cf cor. 1 de Der. 5 (g)}
% archive sheet number: 56
\scanpage{58}
\mspage{64}

% slanted
\margin{``Passage du local au global''}

(\deleted{\unsure{b)}} c) $\Longrightarrow$ a) \deleted{\unsure{Il}} Comme l'hyp. c) est stable par chang$^{\mathrm{mt}}$ de base par imm. ouverte, \deleted{\unsure{on est ramené à}} \inserted{il suffit de} prouver que c) implique que $f$ est une $\unsure{\mathbb D|\underline D}$-équiv, i.e. que pour tout \deleted{\unsure{M}} $F=M_Y$ ($M\in\mathrm{Ob}\,A$), coeff. constant sur $Y$, l'hom d'adj$^{\unsure{\mathrm{on}}}$

\[ M_Y\xrightarrow{\unsure{c}}f_*\underset{M_X}{\underline{f^*M_Y}}\qquad\text{est iso.} \]

Il suffit de \deleted{\unsure{le}} prouver que c'est un iso fibre par fibre (Der 2), donc que pour tt $y\in Y$,

\[ M_y\longrightarrow f_*(M_X)_y \]

est un iso. Par \uline{Der 4 (g)} on a

\[ f_*(M_X)_y\xrightarrow{\sim}\Gamma(X_{/y},M), \]

et l'hyp. que \deleted{\unsure{$f$}} $X_{/y}$ est $\mathbb D$-contractile donne ce qu'on veut.

% slanted
\margin{\uline{\unsure{Prop.} $2'$}}

\uline{\unsure{Dual$^{\mathrm{mt}}$|Dualem$^{\mathrm t}$}}, si $\mathbb D$ est un dérivateur à gauche, conditions équivalentes sur $f:X\longrightarrow Y$

a) $f$ $\mathbb D$-équiv, le reste par tt chgt de base \inserted{$g:Y'\longrightarrow Y$} qui est un \deleted{\unsure{co}} ``isom. \inserted{co}local'' (i.e. $\forall y'\in\mathrm{Ob}\,Y'$, $\underbrace{{}_{y'\backslash}Y'\xrightarrow{\sim}{}_{y\backslash}Y}_{\substack{\text{foncteur induit sur}\\\text{les catégories ``colocalisées''}}}$ (où $y=g(y')$))

b) $\forall y\in Y$, le foncteur induit $ {}_{y\backslash}X\longrightarrow{}_{y\backslash}Y$ est une $\mathbb D$-équiv.

c) $\forall y\in Y$, $ {}_{y\backslash}Y$ est $\mathbb D$-contractile

Pour \unsure{le prouver|la preuve}, on doit utiliser que \fbox{les $ {}_{y\backslash}Y$ sont $\mathbb D$-contractiles.}

\margin{$\swarrow$ Cf. cor. 1 de Dér 5 (d)}
% archive sheet number: 57
\scanpage{59}
\mspage{65}

% slanted
\margin{Meilleure terminologie : \unsure{$f$ est|fonct} $\unsure{\mathbb D|\underline D}$-\uline{asphérique} ($\mathbb D$-équiv stricte), ou $\mathbb D$-\uline{coasphérique} ($\mathbb D$-équiv. co-stricte)}

\uline{Définition 3} \inserted{Soit $\mathbb D$ un prédérivateur (provisoir$^{\mathrm{mt}}$ un dér. à g. ou à dr.)} On dit que $f$ est une $\mathbb D$-équiv. \uline{stricte}, \deleted{\unsure{ou}} \deleted{que $f$ est un foncteur $\mathbb D$-cofinal}, si $f$ satisfait les conditions \inserted{a)} \deleted{équivalentes} de prop. 2. On dit que c'est une $\mathbb D$-équiv costricte si $f$ satisfait celles de prop. $2'$ ; on dit aussi parfois que $f$ est $\mathbb D$-cofinal. (la condition duale, i.e. $\mathbb D$-équiv. stricte, s'appellerait co-$\mathbb D$-cofinal !)

\uline{Prop. 3} \deleted{\unsure{Soit}} \unsure{Soit} $X\longrightarrow Y\longrightarrow Z$. Si $f,g$ sont des $\mathbb D$-équivalences strictes, alors $gf$ aussi. Si $gf$ et $f$ sont des $\mathbb D$-équiv. strictes, alors $g$ aussi.

\inserted{si $\mathbb D$ dér. à g. ou à dr.} \deleted{Toute \deleted{\unsure{mor}} \inserted{équivalence} est une $\mathbb D$-équivalence stricte (en effet $f:X\longrightarrow Y$ une équiv., alors $\forall y\in Y$, $X_{/y}$ est équiv. à la cat. ponctuelle, et on utilise que une \uline{cat. équiv. à la cat. ponctuelle est $\mathbb D$-contractile})}

% slanted
\margin{\unsure{Cas où} $f$ \unsure{a} un adj. à droite.}

% slanted
\margin{$\swarrow$ cf cor. 1 de Dér. 5 (g)}

On utilise le fait que la notion d'isom. local est stable par chgt. de base.

\deleted{\uline{Prop. 4} \inserted{($\mathbb D$ dér. à g.)} Soit $f:X\longrightarrow Y$ une $\mathbb D\unsure{^{+}|^{*}}$-équiv. costricte.}

\diagram{\begin{tikzcd}[ampersand replacement=\&,row sep=tiny]
\deleted{\text{Alors }f^*:\mathbb D(Y)} \arrow[rr,dash] \arrow[dr,dotted,"p"'] \& {} \& \deleted{\mathbb D(X)} \arrow[dl,dotted,"q"] \arrow[r,phantom] \& \deleted{\text{est pl. fidèle}} \\
{} \& \deleted{e} \& {} \& {}
\end{tikzcd}}

\deleted{Alors pour $F\in\mathrm{Ob}\,\mathbb D(Y)$, \inserted{(NB $F$ \uline{\unsure{non|pas} néc. constant})} le morphisme \deleted{\unsure{de}}}

\[ \deleted{q_*(F)\xrightarrow{\sim}p_*(f^*F)\quad\text{i.e. }H^\bullet_{\mathbb D}(Y,F)\xrightarrow{\sim}H^\bullet_{\mathbb D}(X,f^*F)} \]

\deleted{est un iso.}

\deleted{C'est un ``rappel'' de l'axiome Dér 5 (g) pour les dérivateurs à gauche. Si $Z$ est catégorie avec objet final $e$, alors appliquant \uline{Der} 5 \unsure{g)}, à $\{e\}\longrightarrow Z$, on trouve que pour $F$ sur $Z$, on a $H^\bullet_{\mathbb D}(Z,F)\simeq F_e$, ce qui (\unsure{prenant|posant} $F$ constant) implique que $Z$ est $\mathbb D$-contractile.}
% archive sheet number: 58
\scanpage{60}
\mspage{66}

\deleted{Axiome suivant} \deleted{\uline{Der $5'$} (g) : \inserted{$\unsure{\mathbb D|\underline D}$ dér. à g.} Soit $f:X\longrightarrow Y$ tel que $f$ soit}

\uline{Proposition 4} \inserted{$\mathbb D$ dér. à \unsure{g \& d.|dr.|g et dr.}} Soit $f:X\longrightarrow Y$ un foncteur de $\mathrm{Diag}$, tel que pour $y\in Y$, $ {}_{y\backslash}X$ soit $\mathbb D$-contractile (donc $f$ une $\mathbb D$-équivalence costricte, sous réserve que $\mathbb D$ soit dérivateur à \unsure{droit|droite} \ldots) Alors pour \uline{tout} $F$ de $\mathbb D(Y)$ (pas néc. à coeff. constants) on a

\[ H^\bullet_{\mathbb D}(Y,F)\xrightarrow{\sim}H^\bullet_{\mathbb D}(X,f^*(F)) \]

est un iso. \unsure{Réciproquem$^{\mathrm t}$|Réciproquement} si la conclusion est vraie pour tout $F$, alors \deleted{$f$ est une} les $ {}_{y\backslash}X$ sont $\mathbb D$-contractiles.

Dém. \deleted{\unsure{Il}} revient au $=$ de dire que pour tout $M\in\mathrm{Ob}\,A=\mathbb D(e)$, on a

\[ \mathrm{Hom}(M,H^\bullet_{\mathbb D}(Y,F))\xrightarrow{\mathrm{iso}}\mathrm{Hom}(M,H^\bullet_{\mathbb D}(Y,\unsure{f_{\scriptscriptstyle -}|F|\cdot})) \]

\diagram{\begin{tikzcd}[ampersand replacement=\&,row sep=tiny,column sep=small]
X \arrow[rr,"f"] \arrow[dr,dash,"p"'] \& {} \& Y \arrow[dl,dash,"q"] \\
{} \& e \& {}
\end{tikzcd}}

Mais par adjonction, cette flèche s'interprète comme

\[ \mathrm{Hom}(M_Y,F)\longrightarrow\mathrm{Hom}(M_X,f^*F)\tag*{($*$)} \]

D'autre part, \deleted{$M_X=f^*(M_Y)$, \unsure{d'où}}, si $f_!$ existe, on a

\[ \mathrm{Hom}(M_X,f^*F)=\mathrm{Hom}(\deleted{\unsure{M}}f_!(M_X),F), \]

et l'hom ($*$) est déduit de l'hom d'adjonction

\[ f_!(\underbrace{\deleted{\unsure{M}}f^*(M_Y)}_{M_X})\longrightarrow M_Y,\tag*{($**$)} \]

en prenant les Hom dans $F$. Donc la conclusion équivaut à dire que cet hom est iso, quel que soit \deleted{\unsure{$M$}} le \deleted{\unsure{$F$}} coeff constant $M_Y$ sur $Y$. Pour voir si c'est un iso, on regarde fibre par fibre (Der 2), -- or les fibres s'interprètent comme $H_\bullet^{\mathbb D}({}_{y\backslash}X,M_X)$ \inserted{(par \uline{Der 4} (d))} et comme \deleted{\unsure{$y$}}$\unsure{M_{Y,y}|M_y}$.
% archive sheet number: 59
\scanpage{61}
\mspage{67}

Donc la flèche envisagée ($**$) s'interprète comme

\[ H_\bullet^{\unsure{\mathbb D|\underline D}}({}_{y\backslash}X,M_X)\longrightarrow M \]

et c'est l'hom d'adjonction $p_!(p^*M))\unsure{\relbar\joinrel\relbar|\longrightarrow}M$. Donc c'est un iso pour tt $M$ ssi $ {}_{y\backslash}X$ est \deleted{\unsure{$H_{\mathbb D}$}} $\mathbb D$-contractile.

On va \deleted{\unsure{ré}} énoncer ce qu'on a prouvé

\uline{Théorème 1} Soit

\diagram{\begin{tikzcd}[ampersand replacement=\&,row sep=tiny,column sep=small]
f:X \arrow[rr,"f"] \arrow[dr,dash,"p"'] \& {} \& Y \arrow[dl,"q"] \\
{} \& e \& {}
\end{tikzcd}}

dans $\mathrm{Diag}$. On suppose que $\mathbb D$ est un \deleted{\unsure{dér}} prédérivateur satisfaisant Der $(2)_{\unsure{x|X|*}}$ \deleted{et Der \unsure{(4)(d)} pour $f$} que \deleted{\unsure{$f_*$}} $p_*,q_*$ et $f_!$ existent, que $f_!$ satisfait à Der 4 (g) pour le calcul des fibres. Conditions équivalentes

a) $f$ est une $\mathbb D$-équiv. costricte, i.e. (prop. 2) $\forall y\in\mathrm{Ob}\,Y$, $ {}_{y\backslash}X$ est $\mathbb D$-contractile

b) $\forall F\in\mathrm{Ob}\,\mathbb D(Y)$, l'hom suivant est iso

\[ H^\bullet_{\mathbb D}(Y,F)\mathrel{\overset{\sim}{\unsure{\relbar\joinrel\relbar|\longrightarrow}}}H^\bullet_{\mathbb D}(X,f^*F) \]

c) $\forall M\in\mathrm{Ob}(A=\mathbb D(e))$, l'hom d'adjonction est iso.

\[ \alpha_!\underbrace{\alpha^*(M_Y)}_{M_X}\xrightarrow{\sim}M_Y. \]

Je vais renuméroter les \unsure{«|4} axiomes sur les prédérivateurs,

Der 1, Dér 2 comme avant (p. 45)

% slanted
\margin{Axiome des \deleted{diag} cat. d'indices asphériques}

\uline{Dér 3} : Soit $X$ dans $\mathrm{Diag}$, $p:X\mathrel{\relbar\joinrel\relbar}e$. \inserted{Supposons $X$ \uline{contractile}} Alors

\[ p^*:M\longmapsto M_X:\underset{\substack{\Vert\\\mathbb D(e)}}{A}\longrightarrow\mathbb D(X) \]

est pleinement fidèle.

\margin{\deleted{\unsure{En}} \unsure{présence|préciser} \unsure{de|des} Dér (3,4) NB Il suffit de supposer que $X$ a un objet \uline{final}. Mais ce n'est \deleted{\illegible{3}}}

% slanted
\margin{pas un axiome \unsure{autodual|auto-dual} !}

\margin{$\searrow$ cf \uline{variantes} Der a) b) c) page 94, et \unsure{suivantes|surtout} p$^{\mathrm s}$ \deleted{\unsure{96}} \unsure{111--112|111,112}}
% archive sheet number: 60
\scanpage{62}
\mspage{68}

\diagram{\(\begin{array}{lll}
\text{Der 4 (g)}&:&\text{existence des } f_*\\
\text{Der 4 (d)}&:&\text{existence des } f_!\\
\text{\uline{Der} 5 (g)}&&\text{Calcul des fibres de } \deleted{\unsure{R|F}} f_*(F)\\
\text{Der 5 (d)}&&\hspace{1cm}\rule{1.6cm}{0.4pt}\quad \deleted{\unsure{R|F|L}} f_!(F)
\end{array}\)}

L'``ancien axiome Der 5 (g)'' résulte de Der 2, Der 3, Der 4 (g,d), Der 5 (d), grâce au théorème précédent.

Un aspect intéressant du théorème précédent, si on l'envisage comme \unsure{énonçant} l'équivalence de a) et b), c'est que cet énoncé ne fait pas appel à des notions homologiques, mais que sa démonstration y fait appel, sous une forme relativement forte (calcul des fibres de $f_!(F)$). Pour les th. de chang$^{\mathrm t}$ de base pour la $\unsure{\mathbb D|\underline D}$-cohomologie relative i.e. pour les $f_*^{\mathbb D}$), on a besoin aussi constamment du calcul des fibres des $f_*(F)$. Donc il semble qu'on ne peut se passer d'aucun des axiomes précédents, y compris Der 4 (d) et Der 5 (d), même si on ne \inserted{veut} s'intéresser qu'au formalisme des $f^*,f_*$ (aspect cohomologique). Par la suite, on va donc supposer ces axiomes satisfaits.

\uline{Théorème 2} Soit $f:X\unsure{\relbar\joinrel\relbar|\longrightarrow}Y$ dans $\mathrm{Diag}$. Conditions équivalentes

a) $\forall y_0\in\mathrm{Ob}\,Y$, le foncteur canonique \deleted{\unsure{trivial}}

\[ \underset{\text{fibre en }y_0}{X_{y_0}}\xrightarrow{\unsure{\varphi|q}_{\unsure{y|y_0|0}}}X_{/y_0}\qquad (=X\times_Y Y_{/y_0}) \]

est une $\mathbb D$-équivalence costricte.

b) $\forall x_0\in\mathrm{Ob}\,X$, le foncteur induit sur les catégories colocalisées en $x_0$ et en $y_0=f(x_0)$

\[ {}_{x_0\backslash}X\longrightarrow{}_{y_0\backslash}Y \]

est à fibres $\mathbb D$-\unsure{contractiles|contractibles}.
% archive sheet number: 61
\scanpage{63}
\mspage{69}

c) La formation de $f_*F$ \inserted{($F$ \deleted{\unsure{$\in\mathrm{Ob}$}} de $\unsure{\mathbb D|\underline D}(X)$)} commute à tout changement de base $Y'\longrightarrow Y$, et ça reste vrai \deleted{\unsure{pour}} après \deleted{tout} chgt de base -- i.e. si on a $Y''\longrightarrow Y'\longrightarrow Y$ d'où

\diagram{\begin{tikzcd}[ampersand replacement=\&,row sep=tiny,column sep=small]
X'' \arrow[r] \arrow[d,"f''"'] \& X' \arrow[r] \arrow[d,"f'"] \& X \arrow[d,"f"] \\
Y'' \arrow[r,"p'"'] \& Y' \arrow[r,"p"'] \& Y
\end{tikzcd}}

, ``le th. de chgt de base'' est vrai pour $f'_*$ et $p''^*$.

% slanted
\margin{c$'$) La formation de $f_*(F)$ ($F\in\mathrm{Ob}\,\mathbb D(X)$) commute au passage aux fibres, et ça reste vrai après \deleted{tout} chgt. de base $p:Y'\longrightarrow Y$. c$''$) Comme c$'$), avec $Y'\longrightarrow Y$ \unsure{une injection|un morphisme|une ingression} de localisation $Y_{/y}\longrightarrow Y$. d) Si $Y''\xrightarrow{p}Y'\longrightarrow Y$, avec $p$ \unsure{coasph.|coasphér.}, alors $q:X''=X\times_Y Y''\longrightarrow X'=X\times_Y Y'$ \unsure{coasphérique} (\deleted{\unsure{donc}} \unsure{Utilisation|Utiliser} \unsure{du th. 3}). d$'$) Comme d), avec $Y''=e$ (donc) $Y'$ cat. avec objet final $e$ \unsure{$/p$}). \unsure{d$'$)|d$''$)} Comme d), avec $Y''=\unsure{\{i\}|\{1\}|\{e\}}$ et $Y'=\Delta^1$ \unsure{$(p=\{0,1\})$}}

\uline{Dém.} Pour a) $\Longleftrightarrow$ b) cf prop. page 53. Montrons que \deleted{Ces conditions} \deleted{équivalentes} ces conditions impliquent c). Comme elles \unsure{sont} stables par changement de base (cf dém. \deleted{prop.} \inserted{immédiate} \unsure{p.} 55), donc \inserted{par Der 2} on est ramené à prouver \deleted{\unsure{l'iso}} \deleted{que} \deleted{que pour} \unsure{que}

% slanted
\margin{Donc on utilise \uline{Der 5 (d)}}

\[ H^\bullet_{\mathbb D}(X_{/y},F)\mathrel{\overset{\sim}{\unsure{\longleftrightarrow|\longrightarrow}}}H^\bullet_{\mathbb D}(Y_{/y},F)\qquad(F\in\mathrm{Ob}\,\mathbb D(Y)), \]

ce qui résulte du \uline{th. 1}*. (Cf dém. du th. p. \unsure{56|16}, qui marche \unsure{avec ces|avec des} hyp. plus générales.)

On a $c\underset{\vphantom g}{\overset{\text{par Der (2)}}{\Longleftrightarrow}}c'\Longrightarrow c''$

\deleted{Prouvons $c''\Longrightarrow a$. L'hyp. de chgt de base pour $\{y_0\}\longrightarrow Y$, plus \deleted{\unsure{Der}} calcul de la fibre de \deleted{\unsure{$R$}} $f_*(F)$ en $y$ (\uline{Der 5 (g)}), donne que}

\[ \deleted{H^\bullet_{\mathbb D}(X_{/y_0},\deleted{\unsure{F}}G)\longrightarrow H^\bullet_{\mathbb D}(X_{y_0},\varphi^*G)} \]

\deleted{est iso, chaque fois que $G$ est un objet de $\mathbb D(X_{/y_0})$ qui provient d'un $F$ dans $\mathbb D(X)$ par \deleted{\unsure{le foncteur}} \deleted{\unsure{$p$}} $q^*$, où}

\[ \deleted{q:X_{/y_0}\longrightarrow X} \]

\deleted{est le foncteur canonique. Donc on gagne par le th. 1, si on prouve que $q^*$ est ess. surjectif, Il suffit de prouver que pour $G$ sur $X_{/y_0}$, on a}

\[ \deleted{q^*q_*(G)\xrightarrow{\sim}G\qquad\text{(adjonction)}} \]

\deleted{Or pour $(x_0,u_0)\in\mathrm{Ob}(X_{/y_0})$, \inserted{($u_0:f(x_0)\longrightarrow y_0$)} on a, comme $q(x_0,u_0)=x_0$}

\[ \deleted{q^*(q_*(G))_{(x_0,u_0)}=q_*(G)_{x_0}\simeq H^\bullet_{\mathbb D}((X_{/y_0})_{/x_0})} \]
% archive sheet number: 62
\scanpage{64}
\mspage{70}

\margin{\figure{\begin{tikzpicture}[x=0.24cm,y=0.24cm]
\draw[gray,line width=0.7pt,line cap=round] (0.62,1.9) .. controls (0.86,1.61) and (0.84,1.36) .. (0.9,1.2) .. controls (1.28,1.51) and (1.32,1.64) .. (1.39,1.74);
\draw[gray,line width=0.7pt,line cap=round] (0.05,1.57) .. controls (0.51,1.3) and (0.92,1.18) .. (1.4,1.02);
\draw[gray,line width=0.7pt,line cap=round] (0.91,1.22) .. controls (0.98,0.74) and (0.98,0.12) .. (0.53,0.13) .. controls (0.01,0.15) and (0.33,0.75) .. (0.91,1.2);
\end{tikzpicture}}}

\deleted{Appliquons ceci au foncteur déduit de $f:X\to Y$ par le changement de base $Y_{/y_0}\longrightarrow Y$, \unsure{d'où}} \inserted{Considérons} le diagramme à carrés cartésiens

\diagram{\begin{tikzcd}[ampersand replacement=\&]
X_{y_0}\arrow[r,"q"]\arrow[d,"p"] \& X_{/y_0}\arrow[r]\arrow[d,"f_0"] \& X\arrow[d,"f"]\\
\{y_0\}\arrow[r] \& Y_{/y_0}\arrow[r] \& Y,
\end{tikzcd}}

\deleted{de cette façon on trouve} \inserted{L'hyp. c$''$) prouve} que pour tout $F\in\unsure{\mathbb D|\underline D}(X_{/y_0})$ l'on a $H^\bullet_{\mathbb D}(X_{/y_0},F)\xrightarrow{\sim}H^\bullet_{\mathbb D}(X_{y_0},F)$. Par le th. 1, \unsure{cela|celà} implique que $q$ est une $\mathbb D$-équiv. co-stricte.

\uline{Cor. 1} Énoncé dual du th., pour les \unsure{deux|de.} foncteurs $X_{y_0}\longrightarrow{}_{y_0\backslash}X$, $X_{/x_0}\longrightarrow Y_{/y_0}$, et le chang\unsure{emt|ment} de \deleted{\unsure{base}} dans $f_!(F)$ \ldots

\uline{Déf. 4} Sous les conditions \inserted{équiv.} du th. 1, on dit que $f$ est \deleted{\unsure{cohomologiquemt|catérologiquemt}} \uline{$\mathbb D$-propre}. La condition duale s'appelle $f$ $\mathbb D$-lisse.

\deleted{\uline{Prop.} Notons d'abord que chgt de base, par composition}

\uline{Cor. 2} : Soit $f:X\longrightarrow Y$ \unsure{dans|du} $\mathbb D$. Si $f$ est $\mathbb D$-propre \deleted{\unsure{fibré}} ou $\mathbb D$-lisse, et à fibres $\mathbb D$-contractiles, alors $f$ est une $\mathbb D$-équce stricte et costricte. Mieux\unsure{:|,} $f$ est une $\mathbb D$-équce \uline{universelle} i.e. \deleted{\unsure{pour}} c'est une

% slanted
\margin{\uline{Cor 2$'$} \uline{\unsure{généralisé|généraliser}} : Soit \diagram{\begin{tikzcd}[ampersand replacement=\&,column sep=tiny,row sep=tiny]
X\arrow[rr,"f"]\arrow[dr] \& {} \& Y\arrow[dl]\\
{} \& S \& {}
\end{tikzcd}} avec $X,Y$ \unsure{propres|fibrés} sur $S$. Si $f$ induit une $\mathbb D$-équiv. sur les fibres, \unsure{c'est} une $\mathbb D$-équiv. \unsure{Idem} dans le cas lisse. \unsure{C'est le principe|C'est le privilège} \unsure{des conditions|des conclusions}.}

\deleted{Il suffit de voir que $f$ est une $\mathbb D$-équivalence, cela le restera par tt chgt de base \unsure{vu que} les énoncés \unsure{sur} $\mathbb D$-propre et $\mathbb D$-lisse.} $\mathbb D$-équiv, et le reste par \uline{tout} chgt de \inserted{base $Y'\to Y$.}
% archive sheet number: 63
\scanpage{65}
\mspage{71}

\inserted{Il suffit de prouver que $f$ est une $\unsure{\mathbb D|\underline D}$-équiv.} \deleted{\unsure{Ceci s'applique}} \inserted{\unsure{Prenons}} p. ex. $f$ $\mathbb D$-propre. \deleted{Pour ceci,} \inserted{Il suffit} de prouver que si $Y$ est $\mathbb D$-contractile, alors $X$ l'est \deleted{que les $X_{/y_0}$ sont $\mathbb D$-contractiles. Or on} (quitte à l'appliquer \unsure{aux|au} $Y_{/y_0}$ \deleted{\unsure{au lieu de $Y$}}, après chgt de base \inserted{pour voir que les $X_{/y_0}$ sont contr.}).

\diagram{\begin{tikzcd}[ampersand replacement=\&]
\deleted{X_{/y_0}}\arrow[d,phantom,"{\deleted{\downarrow}}" description,"{\deleted{\unsure{f}}}"'] \& X\arrow[d,"f"']\\
\deleted{Y_{/y_0}} \& Y
\end{tikzcd}}

Pour ceci, il suffit de prouver que $f$ est une $\mathbb D$-équivalence, et on sait que ceci est le cas si pour tt \deleted{$F$} \uline{constant} sur

\deleted{\uline{Corollaire 3} Si $f$ est}

$Y$, \unsure{flèche|flèches} $F\xrightarrow{\mathrm{adj}} f_*f^*F$ est iso. En fait, c'est vrai pour \uline{tout} $F$ \unsure{;|:}

\uline{Corollaire 3} : Si $f:X\longrightarrow Y$ est \deleted{\unsure{co}} $\mathbb D$-propre ou $\mathbb D$-lisse, \inserted{à fibres $\mathbb D$-contractiles, alors} $f^*:\mathbb D(Y)\longrightarrow\mathbb D(X)$ est \uline{pl. fidèle}, i.e. les flèches d'adjonction

\[
 F\longrightarrow f_*f^*F,\qquad f_!f^*F\longrightarrow F
\]

dans $\mathbb D(Y)$ sont des iso.

Il suffit de le voir pour $f$ propre, et \unsure{cela|celà} résulte du calcul des fibres de $f_*f^*F$ \ldots

% slanted
\margin{\deleted{\unsure{aux seules fibres|aux fibres}}}

\deleted{\unsure{Prop.}} \uline{Prop. 5} : Les notions de foncteur $\mathbb D$-propre (resp. $\mathbb D$-lisse) est stable par changement de base (trivial), et par composition (utiliser Cor. 2). \unsure{Toute|Tout} \deleted{iso foncteur \illegible{1} fidèle est à la fois $\mathbb D$-propre et $\mathbb D$-lisse} \unsure{imm. fermée|sous-foncteur} est \deleted{\unsure{prop}} $\mathbb D$-propre, \unsure{tt} \unsure{imm. ouv.} $\mathbb D$-lisse. Plus gén, tout iso. local est $\mathbb D$-lisse; \deleted{tout} \unsure{co-}iso. local est $\mathbb D$-propre.

\uline{Prop. 6} Soient $X,Y$ sur $S$, $f:X\to Y$ un $S$-morphisme. Supposons $X$ et $Y$ propres sur $S$ \deleted{\unsure{lisses, et que}} \inserted{p.ex.} ou $X$ et $Y$ lisses sur $S$, et que $\forall s\in S$, $f_s:X_s\longrightarrow Y_s$ soit une $\mathbb D$-équiv, alors $f$ en est une $ {}^*$.

% margin note written sideways
\margin{$*$ \uline{\unsure{Dire|Disons}} : \deleted{Dans le cas $X,Y$ propres sur $S$} \unsure{au principal|au point que} $\equiv$ que sur la sous-catégorie des coeffts sur $Y$ qui proviennent de $\mathbb D(Y)$, $f^*:\mathbb D(Y)\longrightarrow\mathbb D(X)$ induit un foncteur pl. fidèle. Si $f$ propre, pr l'énoncé préc. \unsure{comme|car on a} $q_*q^*F\xrightarrow{\sim}p_*p^*F$ pour tt $F$, et cela résulte du calcul de \unsure{$p_*,q_*$|$q_*,q_*$} par passage aux fibres (\deleted{\unsure{$X,Y$}} \unsure{$q,p$} propres), et de l'hypothèse sur les $f_s$ OK.}
% archive sheet number: 64
